\section{Worked problems}
\subsection*{1. Rounding a tie}
\textbf{Problem.} With scale $0.5$, quantize 1.25 using ties-to-even.
\textbf{Solution.} The unrounded code is 2.5. The even code is 2,
which reconstructs to 1.0. Ties-away-from-zero would instead produce 1.5.
\subsection*{2. Preserve the winner}
\textbf{Problem.} The top-two logit margin is 0.03 and every logit changes
by at most 0.01. Can the winner swap?
\textbf{Solution.} No. The margin can shrink by at most 0.02 and remains
positive. At a margin of exactly 0.02 a tie could occur, so the strict
inequality matters.
\subsection*{3. Distinguish two errors}
\textbf{Problem.} Why does the half-step bound not cover the input 10
with scale 0.5 and codes limited to $[-7,7]$?
\textbf{Solution.} The nearest unbounded code would be 20, but clipping
forces code 7 and reconstruction 3.5. The 6.5 error comes mainly from
clipping, not from rounding to a neighbouring representable code.
\subsection*{4. Decode the scale bytes}
\textbf{Problem.} Interpret little-endian bytes \texttt{00 3E}
as FP16. Why would interpreting the same integer as BF16 be wrong?
\textbf{Solution.} Integer \texttt{0x3E00} has FP16 fields
$s=0,E=15,F=512$, so its value is $2^0(1+512/1024)=1.5$.
BF16 uses different field boundaries and bias. The byte count alone
does not identify the numerical format; BF16 1.5 is \texttt{0x3FC0}.
\subsection*{5. The stopped accumulator}
\textbf{Problem.} Why does adding one repeatedly in FP16 stop at 2,048,
and why does a final FP16 cast of the FP32 result still preserve 4,096?
\textbf{Solution.} At 2,048, spacing is two. The exact 2,049 is a tie,
resolved to the even retained significand at 2,048. Repeating presents
the same tie. FP32 keeps all 4,096 integer partial sums exact; 4,096
itself is a power of two representable in FP16. The lost information
depends on intermediate rounding, not only output type.
\subsection*{6. Bound a projection error}
\textbf{Problem.} For weights $[-1.1,0.2,1.6]$, reconstruction
$[-1,0,1.5]$ and input $[1,2,-1]$, calculate the output error
and its coordinatewise absolute bound.
\textbf{Solution.} The two dot products are $-2.3$ and $-2.5$.
Absolute error is 0.2. The triangle-inequality bound is
$0.1+0.2(2)+0.1=0.6$. Taking absolute values removes cancellation,
so the bound is valid but not attained here. This calculation
excludes arithmetic rounding after reconstruction.
